\begin{table}[htbp!]
\centering
\caption{F1-\textit{score} por categoria de consulta e por modelo}
\label{tab:resultados_por_categoria}
\small
\ajustartabela{%
\begin{tabular}{|l|c|c|c|c|}
\hline
\textbf{Categoria} & \textbf{Qwen 3 4B} & \textbf{Gemma 4 E4B} & \textbf{Gemma 4 E2B} & \textbf{Mistral Nemo 12B} \\
\hline
Simples (n=100) & 0,691 & 0,874 & 0,844 & 0,443 \\
Compostas (n=180) & 0,654 & 0,794 & 0,763 & 0,431 \\
Código MI/INOM (n=57) & 0,949 & 0,878 & 0,929 & 0,362 \\
Tempo relativo (n=60) & 0,725 & 0,547 & 0,645 & 0,450 \\
Ordenação (n=29) & 0,775 & 0,650 & 0,691 & 0,530 \\
Ambíguas/informais (n=197) & 0,640 & 0,825 & 0,744 & 0,429 \\
\hline
\textbf{Média ponderada} & 0,666 & 0,793 & 0,765 & 0,451 \\
\hline
\end{tabular}}
\fonte{Elaborado pelos autores. Uma consulta pode pertencer a mais de uma categoria; n = consultas distintas da categoria nas métricas principais (cada uma pontuada em todas as repetições). As consultas fora do domínio (categoria F), cujo gabarito é não chamar a ferramenta, não têm campos e aparecem na Figura de acurácia por categoria.}
\end{table}
